\subsection{The integral category and finite comparison}

The expression ``integral Betti fibre functor'' requires an integral source. On a nonzero $\Q$-linear category, multiplication by $n$ is an invertible morphism for every positive integer $n$. Its image under an additive functor to free $\Z$-modules would be multiplication by $n$, which is not invertible for $n>1$. Thus no such functor can literally have that source and target. In this section our integral comparisons concern the concrete diagram of smooth projective varieties, integral cohomology modulo torsion, and integral algebraic correspondences, or any integral motivic realization extending it. The counterexample below already occurs on abelian varieties with actual homomorphisms and covariant $H_1(-,\Z)$. Dualizing free lattices gives the identical obstruction for contravariant $H^1$. This states the precise restriction map from a proposed motivic comparison to the diagram we test.

\begin{result}{E.1.1}{Finite diagrams and finitely generated fields}{Proved}
Let $k\subset\C$ have characteristic zero, and let $k_0$ be its algebraic part over $\Q$. Let $v,w:k\hookrightarrow\C$ agree on $k_0$. Every finite diagram of smooth projective varieties over $k$ and algebraic correspondences has an integral Betti comparison between $v$ and $w$, compatible with cup products, K\"unneth maps, duality and those correspondences. If $k/k_0$ is finitely generated, one comparison exists simultaneously for all such varieties and correspondences. It extends to the induced rational tensor category and to the real scalar extensions of the integral realizations.
\end{result}
\begin{proof}
For a finite diagram, choose a finitely generated subfield $F\subset k$ over $k_0$ over which every equation, morphism and correspondence is defined. Because $k_0$ is algebraically closed in $k$, it is algebraically closed in $F$, and characteristic zero makes $F/k_0$ regular. A smooth model $S$ of $F/k_0$ is geometrically integral. After shrinking $S$, the varieties form smooth proper families, the correspondences extend, and all relations used in the finite diagram remain valid. The two embeddings define points of $S(\C)$ after the same embedding of $k_0$; they lie in every nonempty $k_0$-open subset because their maps on $F$ are injective. The complex manifold $S(\C)$ is connected and path connected. Parallel transport in the integral cohomology local systems along a path between those points gives the desired maps. Algebraic correspondences give flat maps of these local systems. Products and the orientation classes of complex fibres are transported, proving all stated compatibilities. Passing to the quotient by torsion and taking duals commute with these maps.

For the simultaneous assertion, take $F=k$ and one smooth model $S/k_0$. The field $k_0$ is countable. List all proper closed $k_0$-subvarieties of $S$; they form a countable family $Z_j$. Each variety and each correspondence under consideration extends over some open subset of $S$, so it is enough to join the two generic points by a path avoiding $\bigcup_j Z_j(\C)$. We give the needed topological argument. In a coordinate ball of complex dimension $d$, a proper complex analytic subset has real Hausdorff dimension at most $2d-2$. For a fixed point outside it, the set of endpoints whose joining segment meets it is contained in a countable union of cones of real Hausdorff dimension at most $2d-1$, and hence has $2d$-dimensional measure zero. This follows by exhausting the analytic subset by bounded pieces of its smooth strata and applying the locally Lipschitz map $(z,t)\mapsto x+t(z-x)$ on bounded intervals of $t$. Countable unions remain null. In each ball we can therefore choose an intermediate point whose two segments avoid all the $Z_j$. A finite chain of overlapping coordinate balls along any initial path, with overlap points chosen outside the same null sets, gives a path avoiding the entire union. Its single parallel-transport system proves the assertion. We use the smooth-proper local-system theorem in the form of Ehresmann's theorem, \cite[Theorem 12.1.1]{Arapura2003}; the argument just supplied is the simultaneous-choice step.
\end{proof}

The finite assertion concerns the specified finite diagram. A category called a ``finite tensor subcategory'' can contain further integral lattices and infinitely many divisibility conditions; one must specify these before identifying it with a finite diagram. The next counterexample is precisely sensitive to those conditions.

\begin{result}{E.1.2}{The universal countable integral assertion}{Refuted}
There is a countable algebraically closed field $k$ of transcendence degree one over $\Qb$, and two distinct complex places agreeing on $\Qb$, for which the integral Betti functors are not naturally isomorphic even on the abelian-variety and isogeny diagram. Each finite subdiagram admits an isomorphism. Thus the unrestricted inference on 15:221 from countability to a global integral isomorphism is false in this explicit interpretation.
\end{result}
\begin{proof}
Take $F=\Qb(t)$ and the elliptic curve
\[
 E/F:\quad y^2=x(x-1)(x-t).
\]
Fix a complex embedding $v$ with $v(t)=e$, and extend it to an algebraic closure $k=\overline F$. We identify the rational constants with the usual $\Qb\subset\C$. The field $k$ is countable, its algebraic part over $\Q$ is exactly $\Qb$, and $v$ is injective because $e$ is transcendental. The curve is smooth since $t(t-1)\ne0$.

Choose an integral homology basis for the Legendre family by transport from a point in $(0,1)$ to $e$. Its local monodromy around $t=0$, in an appropriate oriented basis, is
\[
 T_0=\begin{pmatrix}1&2\\0&1\end{pmatrix}.
\]
The exact integral matrix is recorded in \cite[Example 5.2, pp.~16--17]{OSU2024}. Riemann existence identifies the finite torsion monodromy with the Galois action on the corresponding torsion fields; we use \cite[Corollary 3.7]{OSU2024} and the algebraically closed base-field invariance in \cite[Theorem 1.1]{Landesman2024}, with characteristic zero and the smooth curve $\mathbf P^1-\{0,1,\infty\}$. In particular the image of
\[
 \Gal(F(E_{\mathrm{tors}})/F)\longrightarrow\SL_2(\widehat\Z)
\]
contains the closure of the powers of $T_0$, namely all matrices $u(2a)=\left(\begin{smallmatrix}1&2a\\0&1\end{smallmatrix}\right)$ for $a\in\widehat\Z$. This only uses those actual finite \emph{algebraic} torsion covers; no analytic cover is silently substituted for an algebraic one.

For an explicit noninteger $a$, define its residue modulo $n!$ by
\[
 a\bmod n!=s_n:=\sum_{m=1}^{n-1}m!\pmod {n!}.
\]
The residues are compatible, because $s_{n+1}-s_n=n!$. They define an element of $\widehat\Z$, since the factorial moduli are cofinal among positive integers. For $n\ge4$, $0<s_n<n!/2$, and $s_n\to\infty$. A fixed nonnegative integer has eventually constant least residue; a fixed negative integer has eventually least residue $n!-c>n!/2$. Neither agrees with $s_n$ for all $n$, so $a\notin\Z$. Choose the unique automorphism of the torsion field acting as $u(2a)$ and extend it to $\sigma\in\Gal(k/F)$. Put $w=v\circ\sigma$. These embeddings agree on $F$ but are distinct on the torsion field. They are not complex conjugate, since they both fix $i$.

Every finite cyclic subgroup $H\subset E[n]$ is defined over $k$. Form its algebraic quotient and quotient isogeny
\[
 f_H:E\longrightarrow E/H.
\]
All these varieties and morphisms belong to the concrete diagram. The complex curves $E_v$ and $E_w$ are identical because $v,w$ agree on $F$. Let $\Lambda=H_1(E_v,\Z)=H_1(E_w,\Z)$. Analytic uniformization identifies $f_H$ with the inclusion of lattices
\[
 \Lambda\subset\Lambda_H\subset n^{-1}\Lambda,
 \qquad \Lambda_H/\Lambda=H_v\subset n^{-1}\Lambda/\Lambda.
\]
For $w$ the corresponding subgroup is $H_w=u(2a)H_v$. These statements follow directly by identifying $E_v$ with $\C/\Lambda$ and taking the preimage of $H_v$ in $\C$; the quotient is $\C/\Lambda_H$ and its map is induced by the identity of $\C$.

A natural integral Betti isomorphism would have a component $A\in\GL_2(\Z)$ on $\Lambda$. Naturality for each $f_H$ forces
\[
 A\Lambda_H=\Lambda_{H_w}.
\]
Indeed the component on $E/H$ extends $A$ on the finite-index sublattice $\Lambda$, so its rational extension is forced to be $A$, and integral invertibility gives equality of the lattices. Modulo $n$, this says that $u(2a)^{-1}A$ stabilizes every cyclic subgroup of $(\Z/n)^2$. Stabilization of the three direct summands generated by $e_1,e_2,e_1+e_2$ forces a diagonal matrix with equal unit diagonal entries. Thus
\[
 u(2a)^{-1}A=\lambda I\quad\hbox{in }\GL_2(\widehat\Z),
 \qquad \lambda\in\widehat\Z^{\times}.
\]
Compatibility of the scalars modulo divisors follows from their uniqueness. Consequently $A_{21}=0$ and $A_{11}=A_{22}=\lambda$ are integers. Since $\det A=\pm1$, we obtain $\lambda=\pm1$. Now $A_{12}=\pm2a\in\Z$. An element of $\widehat\Z$ whose double is an integer is an integer: its rational value is integral in $\Z_2$, so that integer is even. Hence $a\in\Z$, the contradiction.

Every finite subdiagram descends to a finite extension of $F$, a finitely generated regular extension of $\Qb$. The two embeddings give points on its connected smooth complex model. E.1.1 supplies the required finite comparison. The finite matrix lemma and the factorial residues are checked by exact arithmetic in \texttt{code/verify\_places.py}; the proof above, rather than any finite bound in that program, establishes the infinite obstruction.
\end{proof}

\begin{result}{E.1.3}{The path-limit step and its completed reading}{Refuted; replacement proved}
Nonempty path sets in the tower on 15:221 need not have a nonempty inverse limit. The completed path torsors do admit compatible points, but those points act on profinite lattices; they need not preserve an integral lattice as E.1.2 demonstrates.
\end{result}
\begin{proof}
For an independent explicit path example set $k'=\bigcup_n\Qb(t^{1/n!})$, with $u_n=t^{1/n!}$ and $u_{n+1}^{n+1}=u_n$. Define
\[
 v(u_n)=\exp(1/n!),\qquad
 w(u_n)=\exp(1/n!+2\pi i s_n/n!).
\]
The residue compatibility makes these field embeddings fixing $\Qb$; their images of $t$ are both $e$. At level $n$, the power map $z\mapsto z^{n!}$ sends a path from $v(u_n)$ to $w(u_n)$ to a loop based at $e$ with winding number congruent to $s_n$ modulo $n!$. A compatible system would therefore have an integer winding number congruent to all $s_n$, impossible by the preceding proof. Each path set is nonempty. Transition maps on these path sets are not surjective: at the first level they miss the disallowed winding classes. The function-field transition has a nontrivial finite algebraic extension, so its generic fibre cannot be geometrically connected; shrinking models cannot change that field extension.

Replacing integer winding numbers by $\widehat\Z$ makes the compatible point $a$ available. More generally the path torsors of a profinite Galois category are nonempty compact spaces. Finite compatibility and compactness supply a point in their inverse limit without a countability requirement. Its scalar extension to $\Z_\ell$ gives the corresponding continuous lattice comparison for each prime. The map from integral transformations to these completed transformations is the completion map; E.1.2 proves that it need not be surjective. This covers both superscript readings $\pi^1$ and $\pi^{\wedge}$ without changing the intended coefficient ring silently.
\end{proof}

\begin{result}{E.1.4}{Uncountable fields: both behaviours occur}{Proved}
There are uncountable fields with distinct complex places inducing isomorphic integral Betti functors on all smooth projective varieties and algebraic correspondences. There are also uncountable fields with two places agreeing on $\Qb$ for which such an isomorphism does not exist.
\end{result}
\begin{proof}
For the positive example choose a transcendence basis $S\cup\{x\}$ of $\R/\Q$, with $S$ uncountable. Let $L$ be the real closure of $\Q(S)$ inside $\R$, put $K=L(i)$, and set $k=K(t)$. The field $K$ is algebraically closed, uncountable, and $x$ is transcendental over $K$. Define $v,w$ to fix $K$ and send $t$ to $x,x+1$, respectively. They are different complex places, both fixing $\Qb$. The polygonal path with vertices
\[
 x,\quad x+ix,\quad x+1+ix,\quad x+1
\]
avoids $K=L+iL$: along either vertical segment the real coordinate is outside $L$, and along the horizontal segment the imaginary coordinate is $x\notin L$. Every smooth projective variety or correspondence over $K(t)$ extends over some open subset of $\A^1_K$; its excluded complex points lie in $K$, since $K$ is algebraically closed. The one displayed path therefore gives parallel transport simultaneously for all of them. Cup products, correspondences and real scalar extensions are compatible as in E.1.1. Neither embedding has real image, so no extra real-place involution is being omitted.

For the negative example, enlarge the algebraically closed countable field $k$ of E.1.2 by an uncountable algebraically independent set and take an algebraic closure. Extend $v$ and $w$ through this extension, using the same independently chosen images for the new transcendence basis and then extending embeddings algebraically. Such choices exist inside $\C$ because its transcendence degree over a countable field is the continuum. Restriction of a natural integral comparison to the original abelian diagram would contradict E.1.2. The algebraic part remains $\Qb$. Thus cardinality alone gives no universal positive or negative assertion.
\end{proof}

\subsection{What the Hodge structure recovers}

\begin{result}{E.3.1}{The labelled integral elliptic test}{Proved}
Let $k$ be any characteristic-zero field admitting complex embeddings. On the diagram of elliptic curves over $k$, the integral Hodge realizations distinguish all \emph{chosen embeddings} $k\hookrightarrow\C$. For a number field, one labelled elliptic object suffices; for a finitely generated field, a finite labelled direct sum suffices. Complex places, defined modulo conjugation, are recovered after making the same quotient of the enriched data.
\end{result}
\begin{proof}
For $a\in k-\{0,1\}$ use the elliptic curve
\[
 E_a:\quad y^2=x^3-\frac{3a}{a-1}x-\frac{2a}{a-1}.
\]
Its discriminant is $1728a^2/(a-1)^3$ and its classical $j$-invariant is $1728a$, as checked symbolically in \texttt{verify\_places.py}. For an elliptic curve $E/\C$, its integral Hodge structure reconstructs its complex torus by the exact quotient
\[
 H^0(E,\Omega^1)^\vee\big/H_1(E,\Z),
 \qquad
 \delta\longmapsto\left(\omega\longmapsto\int_\delta\omega\right).
\]
The Hodge subspace supplies the dual vector space, and the integral pairing supplies the embedded lattice. An integral Hodge isomorphism therefore induces an isomorphism of the complex tori, hence of the elliptic curves, and preserves their $j$-invariants. For embeddings $v,w$, isomorphism of the two Hodge realizations of every $E_a$ forces $1728v(a)=1728w(a)$ for every $a\ne0,1$. They also agree on $0,1$, so $v=w$.

For a number field take a primitive element $a$, adjusting it by a rational integer so that $a\ne0,1$. Its single $E_a$ distinguishes all chosen embeddings. For $k=\Q(a_1,\ldots,a_r)$, apply the same adjustment to each generator and take the direct sum of their $H^1$ objects with the labelled inclusion and projection morphisms. A natural Hodge comparison commutes with each labelled projector and gives the individual isomorphisms; thus it forces equality on every generator. For a conjugate embedding the antiholomorphic comparison exchanges the Hodge labels and conjugates coefficients, exactly as in C.4.1. This identifies the quotient appropriate to unchosen complex places. No passage from integral isomorphism to rational isogeny is made in this proof.
\end{proof}

\begin{result}{E.3.2}{Failure of a universal single object in finite transcendence degree}{Proved}
For $k=\overline{\Qb(t)}$, no motive defined by finitely many algebraic equations, projectors and correspondences distinguishes all chosen complex embeddings through its realization alone. Nevertheless the entire labelled integral Hodge diagram distinguishes them by E.3.1.
\end{result}
\begin{proof}
Any one such object descends to a finite extension $F_0/\Qb(t)$. The algebraic closure $k/F_0$ has a nontrivial automorphism: the function field $F_0$ of a curve over $\Qb$ has, for example, a nonsquare function with a zero of odd order at a smooth point, giving a quadratic extension. Choose a nonidentity $\sigma\in\Gal(k/F_0)$ and set $w=v\circ\sigma$. They give distinct chosen embeddings and identical realizations of the descended object, with its defining finite diagram, because they agree on $F_0$. They fix $i$ and are therefore not related by complex conjugation. This proves the assertion for every individual object. The explicit nonsquare can be obtained from a local uniformizer at a smooth point by the approximation theorem on the curve; its odd valuation proves it cannot be a square.
\end{proof}

\subsection{The proposed injectivity for a large motive}

\begin{result}{E.2.1}{One pair of distinct places defeats every plain integral test}{Refuted; explicit compatible comparison constructed}
There is a countable algebraically closed field $k$ of transcendence degree one over $\Q$, and two chosen embeddings $v,w:k\hookrightarrow\C$ that define different complex places, for which the entire integral Betti correspondence diagrams are tensor-isomorphic. The comparison extends to their real spans and their finite-coefficient comparisons. Thus no single motive, and no collection of plain integral Betti tests on this diagram, separates this pair. This refutes the extension of the injectivity expectation to all fields of finite transcendence degree. The literal 15:223 paragraph assumes a number field; that narrower assertion remains a separate obligation.
\end{result}
\begin{proof}
Put $F=\Qb(t)$ and fix an algebraic closure $k=\overline F$. Start with any embedding $v:k\hookrightarrow\C$ that fixes $\Qb$ and sends $t$ to $e$. On the complex $t$-line take the polygon
\[
 e\longrightarrow e+ie\longrightarrow e+1+ie\longrightarrow e+1.
 \tag{E.2.1}
\]
It avoids $\Qb$: on the vertical pieces its real part is the transcendental number $e$ or $e+1$, and on the horizontal piece its imaginary part is $e$. The real and imaginary parts of an algebraic number are algebraic.

Every element of $k$ is an algebraic function of $t$ over $\Qb$. Its finite branch and pole set in the $t$-line lies in $\Qb$. The value specified by $v$ determines a unique local branch at the transcendental point $e$. Continue this branch along (E.2.1). No branch point or pole is met; its endpoint is finite. Define $w$ to be the endpoint-value map. Addition, multiplication and algebraic identities are preserved by analytic continuation, and a nonzero algebraic function cannot vanish at the transcendental endpoint $e+1$ (its zeros lie over algebraic points). Thus $w$ is an embedding $k\hookrightarrow\C$, fixes $\Qb$, and sends $t$ to $e+1$. Since both embeddings fix $i$, they are not complex conjugates.

A finite collection of varieties and correspondences over $k$ descends to a finite extension $F_0/F$. Normalize a nonempty open subset of the $t$-line in $F_0$, remove its branch points and the finitely many images of bad reduction loci, and spread the collection there. All these removed points belong to $\Qb$. The polygon lifts from the sheet specified by $v$ to the sheet specified by $w$. Proper smooth parallel transport therefore gives integral cohomology isomorphisms, compatible with every correspondence in the finite collection. Enlarging $F_0$ and the collection gives the same maps by uniqueness of path lifting and functoriality of local systems. Hence these maps define one comparison on the whole diagram, including tensor products, duals, projectors for which the lattice is specified, and cup products.

For finite coefficients, Riemann existence identifies transport along the lifted polygon with the corresponding etale path. The definition of $w$ by continuation makes this the same path on every finite cover, so the comparisons commute with the finite-coefficient identifications; inverse limit gives the $\Z_\ell$ statements. Tensoring the integral maps with $\R$ gives the real-span statements. No real-place involution is required, since $k$ contains $i$ and neither embedding has real image. Every motive extracted from this diagram inherits the comparison. The uncountable positive construction in E.1.4 gives a further class of such pairs.
\end{proof}

\begin{result}{E.3.3}{Riemann--Hodge conditions do not characterize geometric places}{Refuted for the explicitly specified enriched target}
Let $\mathcal C$ be the category of pure motivated motives over $\Q$, with its rational Betti fibre functor. On the geometric integral Betti diagram there exists a tensor-compatible assignment of polarizable integral Hodge structures that preserves the usual weight grading, every specified motivic tensor, Tate objects, polarizations, the real Frobenius and the unchanged finite realizations, but does not arise from a complex place of $\Q$. The same construction on the absolute-Hodge abelian subcategory has geometrically realizable real abelian Hodge objects in degree one. Thus the proposed surjectivity fails for the target consisting of these data and the usual Riemann--Hodge relations.
\end{result}
\begin{proof}
Write $G=\Aut^\otimes(\omega_{\B})$ and let $\varepsilon:G\to\Gm$ be its character on $\Q(-1)$. Take
\[
 E:\quad y^2=4x^3-4x-1,\qquad j(E)=110592/37.
\]
The discriminant is $37>0$, and its nonintegral $j$ excludes CM. The group of $H^1(E)$ in the stated category is $\GL_2$; its determinant character is $\varepsilon$. The quotient map $G\to\GL_2$ is faithfully flat by the Tannakian subcategory theorem \cite[Proposition 2.21]{DM1982}. Therefore
\[
 K:=\ker\varepsilon=G\mathbin{\times}_{\GL_2}\SL_2
 \longrightarrow\SL_2
 \tag{E.3.3a}
\]
is faithfully flat.

We justify lifting Lie algebra elements even though $G$ is a proalgebraic group. Over $\Q$ the chosen category admits a countable cofinal sequence of finite tensor subcategories containing $E$ and the Tate object: varieties and algebraic cycles over a number field form countable sets, and every rational tensor space has a countable set of elements. For the absolute-Hodge version the same countability follows from rational Betti tensors. Write the resulting finite quotient groups as $G_n$, and $K_n=\ker\varepsilon_n$. The maps $K_{n+1}\to K_n$ are faithfully flat; in characteristic zero their differentials are surjective. Starting with any $Y\in\mathfrak{sl}_2(\R)$, recursively lift along these surjective real linear maps to obtain $X\in\varprojlim\Lie(K_n)(\R)=\Lie(K)(\R)$.

Let $f=f_\infty\in G(\Q)$ be the natural tensor involution from the real model of every motive over $\Q$. In a rational $f_E$-eigenbasis, the Hodge line of $E$ is spanned by $(1,iy)$ with $y>0$; the real sign condition forces the ratio to be purely imaginary, and nondegeneracy excludes zero. Changing signs of an eigenbasis vector chooses the positive half-plane. Put $\alpha=\sqrt[3]{2}$ and
\[
 a=\sqrt{y/\alpha},\qquad
 Y=\begin{pmatrix}\log a&0\\0&-\log a\end{pmatrix}.
\]
It commutes with $f_E$. Replace the compatible lift $X$ by $(X+\operatorname{Ad}(f)X)/2$. Its projection is still $Y$, it lies in $\Lie K$, and it commutes with $f$. Exponentiating in every finite-dimensional real algebraic group gives compatible elements $g_n=\exp(X_n)$, because algebraic homomorphisms commute with the Lie exponential. They define $g\in K(\R)$ with $gf=fg$ and image $g_E=\operatorname{diag}(a,a^{-1})$.

For each motive $M$, leave its rational and integral Betti groups unchanged and set
\[
 F'^p\omega_{\B}(M)_\C=g_M F^p\omega_{\B}(M)_\C.
 \tag{E.3.3b}
\]
The operator $g_M$ is real, so conjugation exchanges the new Hodge summands just as before. It commutes with every motivic morphism, acts tensor-compatibly, and preserves the central weight grading. The functor is exact because applying an invertible natural transformation preserves strict exact sequences. Since $g\in K$, it acts trivially on every Tate object. For a fixed polarization $q:M\otimes M\to\Q(-n)$, tensor compatibility gives $q(gx,gy)=q(x,y)$. Its new Weil operator is $C'=gCg^{-1}$, and
\[
 (2\pi i)^nq(x,C'x)
 =(2\pi i)^nq(g^{-1}x,Cg^{-1}x)>0\quad(x\ne0).
\]
All polarizations remain positive. The identity $gf=fg$ preserves the anti-Hodge real involution. Integral lattices and finite comparisons have not changed; all their original integral morphisms are Hodge morphisms for the new filtrations by the same naturality argument.

On $E$, equation (E.3.3b) gives ratio $i\alpha$. This is algebraic and not imaginary quadratic: its square is $-\sqrt[3]{4}\notin\Q$, whereas a purely imaginary quadratic number has rational square. Passing from the rational eigenbasis to an integral homology basis is a fractional-linear transformation over $\Q$; it preserves the generated number field and the nonquadratic property. Schneider's theorem (A.3.1) gives a transcendental $j$-invariant for the elliptic torus reconstructed from this integral Hodge structure. Its Riemann form is positive by the calculation above, so it is an elliptic curve over $\C$. An integral Hodge isomorphism to the genuine realization of $E$ would be an isomorphism of elliptic tori and preserve $j$, contradicting the algebraic value $110592/37$. The field $\Q$ has only one chosen complex embedding. Thus this enriched functor does not come from a place.

For degree-one objects of the abelian subcategory, the preserved integral polarizations and real involutions reconstruct abelian varieties over $\R$ by C.1.1; tensor constructions and their Hodge projectors give the corresponding real absolute-Hodge abelian motives. This supplies the stated geometric-object qualification. The fixed algebraic de Rham functor can also be retained, together with a compatible \emph{chosen} comparison: set $\alpha'_M=\alpha_Mg_M^{-1}$. Then $(\alpha'_M)^{-1}F_{\dR}=F'_M$, and $\alpha'$ is a tensor isomorphism preserving the Tate comparison and real involution. It therefore supplies the existence of a comparison satisfying the Riemann--Hodge conditions in the margin of 15:223. It is not asserted to be the distinguished comparison given by integration. Under the forgetful map from enriched functors to plain integral functors this example has the original trivial torsor class; it therefore refutes enriched surjectivity, while it leaves the separate existential question about plain torsor classes open.
\end{proof}

\paragraph{Remaining essential-image problem.}
E.3.1 proves injectivity for the labelled integral Hodge diagram for every field in the prompt. E.2.1 refutes injectivity for the plain integral diagram over general fields of finite transcendence degree; the original number-field assertion is retained separately, and E.3.3 refutes surjectivity for the specified target with the usual Riemann--Hodge conditions, even over $\Q$. These are distinct assertions related by the forgetful morphism from enriched to plain diagrams: E.3.3 retains the original plain functor while changing its filtration. The Riemann--Hodge margin on 15:223 gives no complete axioms for further data that might be intended by ``structure suppl\'ementaire''. The exact image for every such enlarged target, and the image of all integral torsor classes in E.2, remain obligations. We have not replaced that image question by the injectivity proof or by the counterexample to the stated weaker target.

\paragraph{Complex scalar extension and cardinality.}
For a specified neutral Tannakian category, Wibmer's theorem \cite[Theorems 1.1--1.2]{Wibmer2025} proves that two neutral fibre functors over an algebraically closed coefficient field are isomorphic, without a finite-type or countability restriction. Applying it after scalar extension to $\C$ gives the precise comparison map from integral isomorphism torsors to complex ones. The nonexistence in E.1.2 concerns a $\Z$-point; the complexified torsor can have a $\C$-point by this theorem. The theorem therefore closes the complex-coefficient cardinality question and leaves the explicit integral obstruction intact.

\begin{result}{E.3.4}{Including all complex motives does not restore the proposed surjectivity}{Proved counterexample in the specified target}
The deformation of E.3.3 can be constructed on the entire pure motivated category over $\C$, or its absolute-Hodge abelian subcategory, with the same rationally defined elliptic test. It is not induced by any field embedding $\C\hookrightarrow\C$. The construction accommodates the uncountable category without an inverse-limit assumption on integral isomorphisms.
\end{result}
\begin{proof}
Let $G_{\C}$ be the Betti Tannakian group of the chosen category and $K_{\C}$ its Tate-character kernel. The quotient $K_{\C}\to\SL_2$ for the fixed non-CM elliptic curve is faithfully flat. Write its coordinate Hopf algebra as the filtered union of its finitely generated Hopf subalgebras, using \cite[Proposition 2.6 and Corollary 2.7]{DM1982}; take the cofinal subsystem containing the coordinate algebra of $\SL_2$. Every finite-type quotient map to $\SL_2$ is surjective and smooth in characteristic zero. Thus its cotangent map at the identity is injective. Passing to the filtered union gives an injection
\[
 \mathfrak m_{\SL_2}/\mathfrak m_{\SL_2}^2
 \hookrightarrow
 \mathfrak m_{K_{\C,\R}}/\mathfrak m_{K_{\C,\R}}^2.
\]
A real linear functional on the left extends to one on the right by extending a basis. The dual description of Lie algebras therefore lifts the matrix $Y$ of E.3.3 to $X\in\Lie K_{\C}(\R)$, with no cardinality condition. At every finite quotient the matrix exponential is defined, and these exponentials commute with quotient maps, giving a real point of $K_{\C}$.

To retain the real involutions on objects descended from $\R$, let $c^*$ denote conjugation of the base field and let $j_M:\omega_{\B}(M)\to\omega_{\B}(c^*M)$ be the rational tensor comparison induced by conjugating complex points. It satisfies $j_{c^*M}j_M=1$. The formula
\[
 \theta(g)_M=j_M^{-1}g_{c^*M}j_M
\]
defines an involution of $G_{\C}$ preserving $K_{\C}$. On the elliptic quotient it is conjugation by $f_E$. Replace $X$ by $(X+d\theta(X))/2$ before exponentiating. Its image is still $Y$, and the resulting $g$ is fixed by $\theta$. For an object descended from $\R$, naturality with its descent isomorphism identifies this condition with commutation with its real Frobenius.

Now use $F'_M=g_MF_M$ and $\alpha'_M=\alpha_Mg_M^{-1}$ exactly as in E.3.3. The proof there checks tensor compatibility, positivity, integral lattices, finite comparisons, and the chosen de Rham comparison object by object, and therefore applies to this larger category. Every field embedding of $\C$ fixes $\Q$, so its realization of the specified $E/\Q$ has the original algebraic $j$-invariant. The constructed functor gives a transcendental $j$ on this same labelled object. Integral Hodge reconstruction makes an isomorphism to any such geometric functor impossible. This addresses the instruction on 15:227 to include motives over $\C$ in the specified polarizable-Hodge target.
\end{proof}
